\chapter{Compilation of ANN Designs to C}
\label{chapter:appendix_3}
\glsresetall

\section{nn2C Compiler}
Algorithms~\ref{alg:nn2c}, \ref{alg:parse-c}, \ref{alg:write-c}, and \ref{alg:generate-c} present the high-level pseudocode of the \emph{nn2C} compiler \footnote{The compiler can be imported in Python from https://test.pypi.org/simple/nn2c==2.1.6}.  
The tool translates Artificial Neural Network (ANN) designs—single or multi-model—authored in TensorFlow/Keras or PyTorch into synthesizable, fixed-point C code.   An intermediate JSON representation (IR) captures both inter-model wiring and intra-model layer connectivity.\footnote{Supported activations: linear, tanh, sigmoid, exponential/softmax, ELU, ReLU in purely feed-forward MLP or CNN topologies.}

\begin{algorithm}[htbp]
    \caption{Compile TensorFlow/PyTorch Models to C}\label{alg:nn2c}
    \begin{algorithmic}[1]
    \Function{nn2C}{$models$, $conn$, $aParts$, $fname$, $outDir$, $verify$, $simSteps$}
        \State $iRep \gets$ \Call{parse}{$models$, $conn$, $aParts$}
            %\Comment{Build intermediate representation (IR).}
        \State $cFiles \gets$ \Call{writeC}{$iRep$, $fname$, $outDir$, $verify$, $simSteps$}
            %\Comment{Generate all \texttt{.c/.h} sources, Makefile, test benches.}
        \State \Return $cFiles$
            %\Comment{Paths to generated artefacts.}
    \EndFunction
    \end{algorithmic}
\end{algorithm}

The compiler proceeds in two main phases:

\begin{enumerate}
  \item \textbf{Model Parsing} (Alg.~\ref{alg:parse-c}): Builds the IR by extracting architecture, weights and piece-wise linear (PWL) activation tables.
  \item \textbf{C Code Generation} (Algs.~\ref{alg:write-c}-\ref{alg:generate-c}): Emits C/\texttt{.h} sources, a fixed-point arithmetic library, Makefile, and optional SPIN verification harness.
\end{enumerate}

\section{Interface Algorithm}

\begin{algorithm}[htbp]
    \caption{Parse Models to Intermediate Representation}\label{alg:parse-c}
    \begin{algorithmic}[1]
    \Function{parse}{$models$, $conn$, $aParts$}
        \State $iRep \gets \{\}$
        \For{$model, mName$ \textbf{in} $models$}
            \State $arch \gets$ \Call{getArch}{$model$}
            \State $weights \gets$ \Call{getWeights}{$model$}
            \State $acts \gets$ \Call{piecewiseActs}{$model$, $aParts$}
            \State $layers \gets$ \Call{layerGraph}{$model$}
            \State $iRep[mName] \gets$ \{\text{``arch"}: $arch$, \text{``weights"}: $weights$,%
                   \text{``activations"}: $acts$,\text{``lConn"}: $layers$\}
        \EndFor
        \State $iRep[\text{``conn"}] \gets conn$
            %\Comment{Store inter-model wiring.}
        \State \Return $iRep$
    \EndFunction
    \end{algorithmic}
\end{algorithm}

Algorithm~\ref{alg:nn2c} orchestrates the whole flow:

\begin{enumerate}
  \item \textbf{Line 2}: \texttt{parse} converts the supplied models into the IR.
  \item \textbf{Line 3}: \texttt{writeC} materialises every artefact required to build and simulate the generated network in C.
  \item \textbf{Line 4}: The function returns paths to all generated files, enabling downstream integration or packaging.
\end{enumerate}

\begin{algorithm}[htbp]
    \caption{Write C Source Files from Intermediate Representation}\label{alg:write-c}
    \begin{algorithmic}[1]
    \Function{writeC}{$iRep$, $fname$, $outDir$, $verify$, $simSteps$}
        \State $mSrc \gets []$ %\Comment{Per-model \texttt{.c/.h}.}
        \State $actSrc \gets []$ %\Comment{Piece-wise activation kernels.}
        \For{$mName$ \textbf{in} $iRep$ \textbf{excluding} ``conn''}
            \State $mIR \gets iRep[mName]$
            \State $code \gets$ \Call{generateC}{$mIR$}
            \State $act \gets$ \Call{genActC}{$mIR[\text{activations"}]$}
            \State \textbf{append} $code$ \textbf{to} $mSrc$
            \State \textbf{append} $act$ \textbf{to} $actSrc$
        \EndFor
        \State $lib \gets$ \Call{genFixedPointLib}{}
        \State $mk  \gets$ \Call{genMakefile}{$fname$, $verify$}
        \If{$verify$}
            \State $verif \gets$ \Call{genVerification}{$fname$, $simSteps$}
        \EndIf
        \State \Return $[mSrc, actSrc, lib, mk, verif]$
    \EndFunction
    \end{algorithmic}
\end{algorithm}

\begin{algorithm}[htbp]
    \caption{Generate C Artefacts}\label{alg:generate-c}
    \begin{algorithmic}[1]
    \Function{generateC}{$modelIR$}
        \Comment{Emit model-specific \texttt{.c/.h} implementing fixed-point inference.}
        \State \Return $modelCode$
    \EndFunction
    \Function{genActC}{$actPieces$}
        \Comment{Emit lookup-table / piece-wise linear activations.}
        \State \Return $actCode$
    \EndFunction
    \Function{genFixedPointLib}{}
        \Comment{Create common arithmetic library (\texttt{CREATE\_FP}, \texttt{FP\_MULT}, \dots).}
        \State \Return $libCode$
    \EndFunction
    \Function{genMakefile}{$projName$, $verify$}
        \Comment{Makefile with build + optional SPIN verification targets.}
        \State \Return $mkCode$
    \EndFunction
    \Function{genVerification}{$projName$, $simSteps$}
        \Comment{Generate \texttt{main.c} harness and Promela model for SPIN.}
        \State \Return $verifFiles$
    \EndFunction
    \end{algorithmic}
\end{algorithm}

The entry function \texttt{nn2C} expects:

\begin{itemize}
  \item \texttt{models}: list of models or \{\emph{arch, weights}\} pairs.
  \item \texttt{conn}: list of inter-model connections in textual form (e.g., \texttt{``input[0:16]\,$\rightarrow$\,model\_L''}).
  \item \texttt{aParts}: number of PWL segments for activation approximation (default~32).
  \item \texttt{fname}: project base name; reused for the top-level C file and Make targets.
  \item \texttt{outDir}: output directory for generated sources.
  \item \texttt{verify}, \texttt{simSteps}: enable SPIN-based run-time verification and set simulation horizon.
\end{itemize}

\section{Model Parsing Algorithm}

Algorithm~\ref{alg:parse-c} builds the IR as follows:

\begin{enumerate}
  \item \textbf{Line 3}: Initialises an empty container \texttt{iRep}.
  \item \textbf{Lines 4-10}: For each model
        \begin{itemize}
          \item \textbf{Line 5}: \texttt{getArch} captures layer type, dimensions, activation.
          \item \textbf{Line 6}: \texttt{getWeights} serialises parameters (dense weights, convolution kernels, biases).
          \item \textbf{Line 7}: \texttt{piecewiseActs} generates PWL tables—the slopes and intercepts for each activation segment calculated via the analytic derivative at segment mid-points.
          \item \textbf{Line 8}: \texttt{layerGraph} records directed edges between layers.
        \end{itemize}
  \item \textbf{Lines 11-12}: Stores the per-model description in \texttt{iRep}.
  \item \textbf{Lines 13-14}: Adds global \texttt{conn} information for multi-model graphs.
\end{enumerate}

Framework detection is automatic: PyTorch graphs are traversed via \texttt{torch.jit.trace} and ONNX-like nodes; Keras models are parsed layer‐by‐layer.  
The outcome is a unified, back-end-agnostic IR that feeds the code generator.

\section{Compositional Model Parsing}

When multiple sub-models form a composite ANN, the parser:

\begin{enumerate}
  \item Extracts each sub-model independently (Lines 4-10).
  \item Writes inter-model edges (Line 13) so that later stages can route outputs of one model to inputs of another.
  \item Maintains a single fixed-point format across all components, ensuring numerical consistency.
\end{enumerate}

This enables hierarchical designs such as ensembles, split inputs, or late-fusion networks.

\section{C Code Generation Algorithm}

Algorithm~\ref{alg:write-c} transforms the IR into compilable C:

\begin{enumerate}
  \item \textbf{Lines 3-4}: Allocate lists for per-model and activation sources.
  \item \textbf{Lines 5-11}: For every model in the IR
        \begin{itemize}
          \item \textbf{Line 7}: \texttt{generateC} produces a pair of \texttt{.c/.h} files that implement inference with 16-bit fractional fixed-point numbers; neurons are time-multiplexed to save area.
          \item \textbf{Line 8}: \texttt{genActC} emits LUT-based or linear-segment kernels for each activation present.
        \end{itemize}
  \item \textbf{Line 12}: \texttt{genFixedPointLib} supplies arithmetic helpers (\texttt{CREATE\_FP}, \texttt{FP\_MULT}, \texttt{FP\_DIV}, \texttt{argmax}).
  \item \textbf{Line 13}: \texttt{genMakefile} wires everything into a build script; optional SPIN targets are added when \texttt{verify} is true.
  \item \textbf{Lines 14-15}: If verification is requested, \texttt{genVerification} creates a Promela model, harness code, and appropriate assertions (e.g.\ accuracy thresholds).
\end{enumerate}

Algorithm~\ref{alg:generate-c} details the helpers:

\begin{itemize}
  \item \texttt{generateC}: Emits state structs for inputs, hidden and output neurons, plus two procedures—\texttt{\emph{Model}Init} (zeroing state) and \texttt{\emph{Model}Run} (one simulation step).
  \item \texttt{genActC}: Converts each PWL table to an \texttt{int32\_t} fixed-point function with branch-less segments where possible.
  \item \texttt{genMakefile}: Provides rules for library compilation, object generation, linking, and cleaning.
  \item \texttt{genVerification}: Builds a C test-bench reading \texttt{test/input.txt} and a Promela script automatically sweeping an input noise range, asserting the desired number of correct predictions.
\end{itemize}

The generated project folder therefore contains:

\begin{itemize}
  \item Top-level \texttt{main.c} (or \texttt{verify/main.c}) calling the model.
  \item Model-specific \texttt{.c/.h} files (\emph{one per sub-model} if compositional).
  \item Activation kernels (\texttt{relu.c}, \texttt{tanh.c}, \dots).
  \item Fixed-point arithmetic library (\texttt{library.c/h}).
  \item Optional \texttt{result.c/h} for classification post-processing.
  \item Makefile with build and simulation targets.
\end{itemize}

The C code is generated using the synchronous approach, as described in Section \ref{sec:hardware}. As this code does not require minimisation of hardware resources, we do not use the multicyclic approach. Instead, we use the pipelined layer-by-layer approach, where each layer is executed at a time in a pipeline. For example, if a model has 5 layers, the pipeline will take 5 ticks to produce the first output.

A minimal usage example is shown in Listing~\ref{lst:usage_example_c}.

\begin{lstlisting}[language=Python, caption=Python invocation of the nn2C compiler, label=lst:usage_example_c]
from nn2C.nnParser import ToC
from tensorflow.keras import layers, Model

def build_mlp():
    x_in  = layers.Input(shape=(16,))
    x     = layers.Dense(16, activation='relu')(x_in)
    x     = layers.Dense(8,  activation='relu')(x)
    x_out = layers.Dense(4,  activation='softmax')(x)
    return Model(x_in, x_out)

net1 = build_mlp()
net2 = build_mlp()

ToC(models=[net1, net2],
    model_names=["net_L", "net_R"],
    connections=["input[0:16] -> net_L", "input[6:22] -> net_R"],
    classes=[0,1,2],             # for optional result block
    approx_parts=32,             # PWL granularity
    fname="ensemble",            # base project name
    result_block=True,
    verify=True,                 # generate SPIN harness
    simulation_steps=3)
\end{lstlisting}
